\section{Persevere in the Mercy You Have Received}

St.~Alphonsus addresses the person who has renounced sin and hopes to have received pardon, then reminds him that the race is not yet over. Perseverance is needed until the end (\emph{Preparation for Death}, XXXI.1). The mercy received should become the manner of life in which the Lord will find us. The resolution formed during recollection has to survive the ordinary return of fatigue, temptation, irritation, and desire.

Alphonsus names avoidance of occasions, frequent recourse to the sacraments, daily meditation, and prayer. He warns against placing confidence simply in one's purposes. We need God's assistance in the conflict, and the discovery of our weakness should make recourse to him more prompt. A person may sincerely resolve to change while still relying upon the same unguarded arrangements in which he repeatedly fell.

An occasion of sin must be judged by its effect upon this person's conduct. An arrangement may repeatedly place him where he is likely to consent to grave evil. He may already know from several falls that a general resolution has left the danger undisturbed. Amendment then requires changing the arrangement, with the inconvenience or loss that follows. The relevant question is what would actually interrupt the familiar course toward consent. Christ's teaching about the hand and the eye gives a reason to accept a painful renunciation for the good of eternal life.

The safeguard should address the temptation before it becomes acute. If a recurring exchange reliably ends in cruel speech, the person can decide beforehand how to leave it when its direction becomes clear. Waiting until anger has taken command makes the refusal harder. This precaution acknowledges what experience has shown about his weakness. It also gives his purpose of amendment a definite act to perform. He must still ask for grace when the provocation comes, but he has ceased deliberately placing himself where he repeatedly abandons the restraint he intended to keep.

Prayer accompanies the safeguard. It acknowledges that the soul cannot produce supernatural fidelity from its own resources. Alphonsus urges immediate recourse to Jesus and to the intercession of his Mother in temptation. The person who has learned this dependence does not need to wait until a fall to remember God. The first movement of danger can become a summons to ask for help, and the help received can become an occasion of gratitude.

Daily examination helps reveal whether the annual resolution has survived. A familiar excuse may have returned before we notice that the old conduct is again becoming ordinary. A brief review of the day can identify the occasion on which we yielded and the help we failed to use. It can also recognize fidelity which discouragement overlooks. The review should lead to sorrow for the fault, thanksgiving for grace, and an act of amendment appropriate to what happened. Its value lies in returning promptly to God, before delay acquires another week's force.

At the yearly recollection, set aside enough quiet to consider one of Christ's warnings. The rich man and Lazarus, for example, require attention to a need which had remained near the rich man's gate. We can ask whether a known claim upon our charity has likewise become easy to pass by. In prayer we can consider how indifference became settled and what responding would now cost. That examination must remain honest about our responsibilities and means. It should also disturb an excuse which has allowed a real duty to remain indefinitely outside our concern.

Consider the finality Christ describes before deciding upon the response. The rich man's request for relief cannot undo the separation, and his concern for his brothers turns toward those who can still hear Moses and the prophets. We are among the living who have received the warning. A duty postponed during recollection will soon face the same competing attractions as before. The thought of the end helps us judge whether those attractions deserve the obedience we have given them. Friendship with God is endangered by a grave refusal, even when the immediate inconvenience of obeying fills our attention.

Bring what has become clear into repentance. If conscious of mortal sin, turn to God in contrition and seek sacramental confession promptly. Let the examination be honest and proportionate; let the accusation be candid; let the purpose of amendment include the needed change. Where restitution or reconciliation is due, determine the next just act and undertake it. If a matter is genuinely difficult, bring it to the confessor's counsel with a willingness to follow sound direction.

Remember the mercy by which a return is possible. A previous confession may have ended a refusal which once seemed fixed. Recalling it gives reason to seek the same Savior's help in the fault now exposed. This memory must include what made the return possible: grace received, the truth spoken, and a willingness to surrender the sin. It can strengthen hope without making the present act of repentance unnecessary. The Lord who received us then is the Lord to whom we must now bring what we have continued to withhold.

Take one definite resolution from the recollection into the coming day and ask for God's help to keep it. Choose the act which the examination has shown to be required, then provide for the circumstances in which it must be done. A resolution to repair a wrong needs a first step and a time; an occasion to be removed needs an actual change. Return to the purpose when its cost becomes apparent. Alphonsus's counsel on perseverance concerns that continuance, when the initial impression has faded and the ordinary duty still asks for consent.

In the prayer following his consideration on abused mercy, St.~Alphonsus asks for the love and perseverance needed to remain faithful. We may make our own this complete sentence from that received prayer:

\begin{quote}
``My God, for the love of Jesus Christ, give me holy perseverance till death, along with your holy love.''
\end{quote}

\noindent\emph{Preparation for Death}, XVII.1, prayer excerpt; anonymous English translation, Dublin, 1843, p.~209.

\clearpage
